% Secao global. Um paragrafo: por que ordenacao importa, o que foi feito, o
% que foi encontrado. Revisar a cada entrega, quando entrar algoritmo novo.

\indent \qquad Ordenar é uma das operações mais frequentes em computação, e os
algoritmos que a realizam diferem muito em custo conforme o tamanho e a ordem
inicial dos dados. Conhecer esse comportamento é o que permite escolher o
algoritmo certo para cada situação. Este trabalho implementa o Insertion Sort
em linguagem C e mede o seu tempo de execução em dezoito combinações de entrada:
seis tamanhos, de 10 a 1.000.000 de elementos, em três tipos de instância ---
crescente, decrescente e aleatória. Os resultados confirmaram a análise teórica.
Com a entrada já ordenada, o algoritmo é linear e não registrou tempo mensurável
em nenhum tamanho. Com entradas decrescente e aleatória, ele é quadrático:
multiplicar o número de elementos por dez multiplicou o tempo por cerca de cem,
e a entrada decrescente, que é o pior caso, custou o dobro da aleatória. O
Insertion Sort mostrou-se adequado a conjuntos pequenos ou quase ordenados e
impraticável para entradas grandes em ordem qualquer, em que um milhão de
elementos levou até dois minutos para ser ordenado.
